\subsection{Norm-compatible cyclotomic units detect the quadratic companions}
\label{sec:universal-cyclotomic-pair}

Odd Iwasawa torsion creates a distinction between arbitrary odd
cohomology classes and the cyclotomic classes used by the modular-symbol
theorem.  The latter annihilate that torsion on restriction to the
cyclotomic kernel.  This supplies the quadratic pole condition under
precise hypotheses of Wake--Wang-Erickson, without assuming that the
whole plus Iwasawa module is finite.  The cubic pole condition remains
a separate arithmetic hypothesis.

\begin{lemma}[K17-NORM-TORSION]\label{lem:norm-torsion}
Let \(P=H\rtimes\Gamma\) be a pro-\(p\) group, where
\(\Gamma=\langle\gamma\rangle\cong\mathbb Z_p\).  Let the coefficient
character \(\kappa\) be trivial on \(H\), and write
\(\kappa(\gamma)=u\in1+p\mathbb Z_p\).
Put \(P_n=H\rtimes\Gamma^{p^n}\),
\(\Lambda=\mathbb Z_p[[\Gamma]]\), and \(M=H^{\rm ab}\).
Suppose that classes
\[
 c_{n,r}\in H^1(P_n,\mathbb Z/p^r(1)),
 \qquad n\geq0,\quad r\geq1,
\]
are compatible with corestriction in \(n\) and coefficient reduction
in \(r\).  Then every restriction
\[
 c_{n,r}|_H:M\longrightarrow\mathbb Z/p^r(1)
\]
annihilates the entire \(\Lambda\)-torsion submodule of \(M\).
\end{lemma}

\begin{proof}
Write \(a_{n,r}\) for this restriction.  Since \(H\) acts trivially
on the coefficients, it is an actual continuous homomorphism and is
unchanged by replacing a cocycle by a cohomologous one.  The cocycle
identity gives
\[
 a_{n,r}(\gamma^{p^n}m)=u^{p^n}a_{n,r}(m).
\]
Choose a compatible Tate basis to write scalar coordinates, and define
\[
 \Phi_{n,r}(m)=
 \sum_{j=0}^{p^n-1}u^{-j}a_{n,r}(\gamma^j m)[\gamma^j]
 \ \in(\mathbb Z/p^r)[\Gamma/\Gamma^{p^n}].                 \tag{NC.1}
\]
Give this coefficient module the \(P\)-action with \(H\) trivial and
\(\gamma\) acting by multiplication by \(u[\gamma]^{-1}\).
Shifting the sum, including its endpoint, proves
\[
 \Phi_{n,r}(\gamma m)=u[\gamma]^{-1}\Phi_{n,r}(m).
                                                               \tag{NC.2}
\]
The restricted corestriction formula, with representatives
\(\gamma^{-kp^n}\), is
\[
 a_{n,r}(m)=
 \sum_{k=0}^{p-1}u^{-kp^n}a_{n+1,r}(\gamma^{kp^n}m).
                                                               \tag{NC.3}
\]
It follows that ordinary group-ring projection from level \(n+1\)
to level \(n\), which sums coefficients over each fiber, sends
\(\Phi_{n+1,r}\) to \(\Phi_{n,r}\).  Coefficient reduction is also
compatible.  The inverse limit of these actual homomorphisms is
therefore a continuous equivariant map
\[
                  \Phi:M\longrightarrow\Lambda^\#(1).     \tag{NC.4}
\]
This is integral Shapiro in restriction coordinates.  The inverse
group action and corestriction convention agree with
\cite[\S2.5, pp.~16--17]{SharifiReciprocity2017};
the finite formulas above prove the part needed here directly.
In particular, no interchange of cohomology and inverse limit or
vanishing of a \(\varprojlim^1\) term is used.

If \(f(\gamma)t=0\), where \(f\in\Lambda\) is nonzero, then
equivariance implies
\[
                   f(u\gamma^{-1})\Phi(t)=0.
\]
The substitution \(\gamma\mapsto u\gamma^{-1}\) is an involutive
continuous automorphism of \(\Lambda\): in the parameter
\(Z=\gamma-1\), it is \(Z\mapsto u/(1+Z)-1\).
Since \(\mathbb Z_p[[Z]]\) is a domain, \(\Phi(t)=0\).
The coefficient of the identity in (NC.1) is \(a_{n,r}(t)\).
Hence every \(a_{n,r}(t)\) is zero, as claimed.
\end{proof}

\paragraph{All finite cyclotomic inputs.}
For \(p\) odd, put \(F^{(n)}=\mathbb Q(\mu_{p^{n+1}})\) and choose
compatible roots of unity.  For every \(a\in\mathbb Z_p^\times\),
\[
 v_n(a)=1-\zeta_{p^{n+1}}^{\,a\bmod p^{n+1}}
 \quad\hbox{satisfies}\quad
 N_{F^{(n+1)}/F^{(n)}}v_{n+1}(a)=v_n(a).                    \tag{NC.5}
\]
Kummer maps commute with norm and corestriction.  Their finite
classes therefore satisfy Lemma~\ref{lem:norm-torsion}.
They are classes of the maximal pro-\(p\) quotient: the Tate
action factors through it, and the degree-one comparison in
Lemma~\ref{lem:etale-prop} is an isomorphism, natural for the
tower transfers.  Equivalently, the Galois closure of each Kummer
extension is a \(p\)-extension, since the base contains \(\mu_p\).

Nonunit exponents cause no loss.  Every nonzero exponent modulo
\(p^{n+1}\) has the form \(p^s b\), with \(0\leq s\leq n\) and
\(b\) a unit, and
\[
 1-\zeta_{p^{n+1}}^{\,p^s b}
  =\prod_{j=0}^{p^s-1}
       \left(1-\zeta_{p^{n+1}}^{\,b+jp^{n+1-s}}\right).
                                                               \tag{NC.6}
\]
Every factor has a unit exponent and is a value of a family in
(NC.5).  Formula (NC.6) is the identity
\(\prod_{j=0}^{p^s-1}(1-X\eta^j)=1-X^{p^s}\) for a primitive
\(p^s\)-th root \(\eta\); it takes place in the same finite field
and involves no division by the degree.  Thus the restriction of
every finite cyclotomic Kummer input to \(H\) annihilates all of
\(T=M_{\Lambda\text{-tors}}\).

Complex conjugation \(c\) commutes with the norms, so averaging
these families by \((1+c)/2\) preserves this property.
For a \(p\)-unit \(v\), the quotient \(v/c(v)\) is a global unit
of absolute value one at every embedding: the valuations at the
unique prime above \(p\) cancel.  It is therefore a root of unity.
The negative parts of the Kummer classes consequently lie in one
cyclic subgroup.  Their mutual cups vanish, because self-cup is
zero when \(2\) is invertible.  Mixed-parity cups lie in the odd
target.  The even target spanned by cyclotomic symbols is thus
spanned by cups of \(c\)-even classes whose restrictions factor
through \(Q=M/T\).

\begin{proposition}[K17-UNIVERSAL-PAIR]\label{prop:universal-pair}
Let \(p\) be odd and \(P=H\rtimes\Gamma\) as above, with \(H\)
free pro-\(p\).  Suppose that
\(\Delta\cong\mathbb F_p^\times\) acts on \(H\), commutes with
\(\Gamma\), and extends \(\kappa\) to \(P\rtimes\Delta\).
Let \(c\) be its involution, and assume \(\kappa(c)=-1\).
Assume that \(M\) is finitely generated over \(\Lambda\), that
\(T=M_{\Lambda\text{-tors}}\) is finite free over \(\mathbb Z_p\),
and that \(Q=M/T\) is torsion free over \(\mathbb Z_p\), has only
negative \(c\)-parity, and is free of finite rank over
\(R=\Lambda[1/p]\) after inverting \(p\).
Assume that every even character component \(T_e^+[1/p]\) is
cyclic over \(R\).

Suppose that, on a cofinal diagonal of group and coefficient
levels,
\[
 H^2(P_n,\mathbb Z/p^n(2))^+
\]
is generated by cups of even classes whose restrictions to \(H\)
annihilate \(T\).  Then the actual free-\(Q\) quadratic companions
satisfy Pair-NV: for every root \(\alpha\) of the characteristic
and minimal polynomial \(q_e\) of \(\gamma\) on
\(T_e^+[1/p]\), at least one matching coefficient of \(C_e\) is
nonzero on
\[
                     (1+T_1)(1+T_2)=\alpha.                \tag{NC.7}
\]
Repeated roots are allowed.  The coefficients are in the
bounded-denominator coordinates of an \(R\)-basis of \(Q[1/p]\).
\end{proposition}

\begin{proof}
\par\medskip\noindent\textbf{The actual extension.}
The class-two quotient of \(H\) has its integral
Baker--Campbell--Hausdorff Lie structure, since \(p\) is odd.
Its central commutator module is
\(\widehat{\bigwedge}^2_{\mathbb Z_p}M\).
Quotient this module by all brackets with a \(T\) input.
The remaining central module is
\[
                      N=\widehat{\bigwedge}^2_{\mathbb Z_p}Q.
\]
For example, exactness follows from a continuous
\(\mathbb Z_p\)-splitting of \(M\to Q\), which exists because
compact torsion-free \(\mathbb Z_p\)-modules are products of
copies of \(\mathbb Z_p\).  No equivariant splitting is asserted.

The central module \(N\) has positive \(c\)-parity.  Hence the
negative part of the quotient Lie module maps isomorphically
onto \(M^-\), giving a canonical \(\Gamma\)-equivariant section.
The lift of \(T^-\) is central, because brackets with \(T\)
were killed.  Quotient this central equivariant lift as well.
The odd part is now the canonical section from \(Q\), and the
even part is the extension
\[
          0\longrightarrow N\longrightarrow E
                 \longrightarrow T^+\longrightarrow0.     \tag{NC.8}
\]
For a cyclic vector \(v_e\) and an even lift \(y_e\) in its
character component, set
\[
                    q_e(\gamma)y_e=C_e\in N[1/p].           \tag{NC.9}
\]
This is the quadratic free-\(Q\) companion used in
Theorem~\ref{thm:odd-torsion-cubic}.  Brackets with torsion
lifts vanish in this quotient, so its ordered companion product
is the additive characteristic-polynomial expression.
The logarithm and Magnus quadratic coefficient agree on a
derived element.

\par\medskip\noindent\textbf{A vanishing residue gives a split socle.}
Pass to a finite splitting field \(L/\mathbb Q_p\), with integers
\(\mathcal O\).  Finite generation clears the denominators of the
coordinates \(Q[1/p]\cong R^d\) and their inverse.  Completed
tensor and exterior squares retain a fixed denominator bound.
Thus \(N[1/p]\) is the alternating bounded-denominator
two-variable coefficient module.  Continuity of the pro-\(p\)
action on \(T\) gives \(\alpha-1\in\mathfrak m_{\mathcal O}\).

Suppose every coefficient of \(C_e\) vanishes on (NC.7).
The integral formal-coordinate change
\[
 S=(1+T_1)(1+T_2)-\alpha,\qquad
 T_2=\frac{\alpha+S}{1+T_1}-1
\]
has specialization kernel \((S)\).  Division preserves the
common denominator bound and, since \(S\) is symmetric, the
alternating property.  Hence
\[
                        C_e=(\gamma-\alpha)A
                        \quad\hbox{in }N_L.
\]
Let \(r_e(z)=q_e(z)/(z-\alpha)\).  Cyclicity gives a nonzero
vector \(w=r_e(\gamma)v_e\) spanning the \(\alpha\)-socle,
including when that root is repeated.  Then
\[
                       \widetilde w=r_e(\gamma)y_e-A
\]
is an equivariant lift of \(w\).  Put \(W=Lw\) and
\(W_{\mathcal O}=W\cap T_{\mathcal O}^+\).
This lattice is saturated, hence a direct summand over
\(\mathcal O\).  Choose an \(\mathcal O\)-linear section
\(s:T_{\mathcal O}^+\to E_L\) extending the equivariant section
on \(W_{\mathcal O}\).  One fixed power \(p^b\) clears its
denominators.  Write the class-two logarithm of \(h\in H\) as
\[
 x(h)+s(t(h))+z(h),\qquad
 x(h)\in Q_{\mathcal O},\quad
 t(h)\in T_{\mathcal O}^+,\quad z(h)\in N_L.
\]
The same \(p^b\) makes every \(z(h)\) integral.  The section
defect
\[
            \beta_\delta(t)=\delta s(t)-s(\delta t)
            \quad(\delta\in\Gamma)
\]
vanishes on \(W\).

\par\medskip\noindent\textbf{The actual cup and its dual quotient.}
Average two even cocycles \(f,g\) under \(c\).
They vanish on the chosen \(\Gamma^{p^n}\) section, since \(c\)
fixes its elements and acts by \(-1\) on the coefficients.
Their restrictions to \(H\) factor through \(Q\), by hypothesis.
Let
\[
 \psi=f\wedge g:N_{\mathcal O}\longrightarrow
                  \mathcal O/p^n(2).
\]
The integral cochain
\[
       k(h)=-\frac{p^b}{2}f(h)g(h)-\psi(p^b z(h))            \tag{NC.10}
\]
has differential \(dk(h_1,h_2)=p^b f(h_1)g(h_2)\).
Indeed the central coordinate of the product is
\(z_1+z_2+[x_1,x_2]/2\).  The two contributions to the
differential are
\[
 \frac{p^b}{2}(f_1g_2+f_2g_1),\qquad
 \frac{p^b}{2}(f_1g_2-f_2g_1).
\]
Extending \(k\) by \(k(h\delta)=k(h)\), its equivariance
defect computes the Hochschild--Serre edge of the cup.
Up to the convention for its sign, this is
\[
                         \psi(p^b\beta_\delta(t)),          \tag{NC.11}
\]
which vanishes on \(W\).

Both \(H\) and \(\Gamma^{p^n}\) have \(p\)-cohomological
dimension at most one.  All odd components of \(M\), including \(T^-\),
give negative parity in \(\operatorname{Hom}(M,\mathcal O/p^n)(2)\).
Therefore
\[
 H^2(P_n,\mathcal O/p^n(2))^+
 \cong H^1\bigl(\Gamma^{p^n},
         \operatorname{Hom}(T^+,\mathcal O/p^n)(2)\bigr).
                                                               \tag{NC.12}
\]
Set \(V=\operatorname{Hom}_{\mathbb Z_p}(T^+,\mathcal O)(2)\),
and let \(B\) denote its \(\gamma\) operator.  Normalized
generator evaluation identifies the last group with
\[
                   V/\bigl(p^nV+(B^{p^n}-1)V\bigr).         \tag{NC.13}
\]
Corestriction sends a generator value to the same value in
the larger group's coinvariants.  Its transition is consequently
the quotient map, with no extra factor \(p\).
As \(B^{p^n}-1\to0\), the denominator submodules are cofinal
with the \(p\)-adic neighborhoods, and the inverse limit is \(V\).

Restriction to \(W_{\mathcal O}\) gives a nonzero torsion-free
quotient
\[
 V\longrightarrow
       \operatorname{Hom}_{\mathcal O}(W_{\mathcal O},\mathcal O)(2).
\]
It is surjective because \(W_{\mathcal O}\) is saturated.
Its rational root is \(u^2/\alpha\), and its character is
\(\kappa^2 e^{-1}\).
By (NC.11) and finite-level cup generation, \(p^b\) times every
finite target has zero image in the corresponding quotient of
this dual line.  The separated inverse limit would make \(p^b\)
annihilate the displayed nonzero torsion-free quotient,
a contradiction.  The same \(b\) works at all levels; no
compatible expressions as sums of cups are required.
This proves (NC.7) for all roots.
\end{proof}

\begin{corollary}[K17-WWE-PAIR]\label{cor:wwe-pair}
Let \(p\geq5\), and use the arithmetic cyclotomic group and
notation of the preceding kernel theorems.
For every even character \(\theta\ne1,\omega^2\), put
\(\chi=\omega^{-1}\theta\).
Let \(\xi_\chi\) be a characteristic series of \(X_\chi(1)\)
in the Wake--Wang-Erickson ring and twist convention.
For every nonzero component suppose that
\[
 \begin{gathered}
 X_\chi(1)[1/p]\text{ is cyclic over }
                         (\Lambda/\xi_\chi)[1/p],\\
 X_\theta\otimes_\Lambda(\Lambda/\xi_\chi)[1/p]=0.
 \end{gathered}                                             \tag{NC.14}
\]
Then every rational even torsion character is cyclic, and the
actual free-\(Q\) quadratic companions satisfy Pair-NV.
This conclusion allows nonzero odd \(\Lambda\)-torsion in \(M\).
\end{corollary}

\begin{proof}
The conditions are exactly those of
\cite[Corollary~8.1.3(2), p.~49]{WakeWangEricksonPseudo2018}.
Its conclusion is an integral isomorphism
\[
 \varpi:
 (\mathcal H^-_\theta/I_\theta\mathcal H^-_\theta)/
           (p\text{-primary torsion})
       \ \xrightarrow{\ \sim\ }\ X_\chi(1).
\]
In particular, the original \(\varpi\) is surjective.
The full proof allows repeated roots; it does not require the
derivative of the characteristic series to be invertible.
At tame level one the primitive-character and other setup
conditions exclude exactly \(1,\omega^2\) among the even
characters \cite[\S1.1, pp.~2--3]{WakeWangEricksonPseudo2018}.
We apply its nonunit-series running hypothesis only to nonzero
components.  A unit characteristic series makes \(X_\chi\)
finite, and its \(p\)-torsion freeness makes it zero.

Fukaya--Kato turn this surjectivity into finite-level generation
by \(\{1-\zeta^a,1-\zeta^b\}\) for every cyclotomic level
\cite[\S\S7.1.8--7.1.10, p.~86]{FukayaKato2012}.
Their exponents are nonzero and at least one is prime to \(p\);
(NC.6) handles every input, including a nonunit exponent.
Their \S7.1.9 identifies the reduction modulo \(p^r\) with
the corresponding class-group quotient.

The excluded even targets \(1,\omega^2\) correspond to
\(X_{\omega^{-1}},X_\omega\), which vanish by the
Stickelberger exception calculation; the first is explicitly
recorded in \cite[proof of Lemma~5.2.7, p.~13]{WakeWangErickson2017}.
Norm projections from \(X^-\) onto every finite minus class
group are surjective
\cite[Proposition~3.4.1 and proof, pp.~90--91]{SharifiIwasawaNotes}.
Thus their finite targets also vanish.  At \(F^{(n)}\), the
Fukaya--Kato coefficient index is \(n+1\); reduction gives
coefficients \(p^n\), as required in
Proposition~\ref{prop:universal-pair}.  The Brauer term is zero,
since there is one omitted finite prime and no real place.

By Lemma~\ref{lem:etale-prop}, Kummer classes lift uniquely to
group \(H^1\), and group \(H^2\) injects into etale \(H^2\).
Lifting the generating cup expressions makes the latter
injection surjective on the even part.  This comparison is
natural for the tower transfers and assumes no arithmetic
\(K(\pi,1)\) theorem.  The paragraph following
Lemma~\ref{lem:norm-torsion} supplies even generating inputs
whose restrictions annihilate \(T\).

The unconditional module calculations in the preceding
arithmetic sections give \(Q[1/p]\) rank one in every odd
character, and zero in even characters.  The first condition
in (NC.14), and the Kummer pseudoisomorphism
\[
                         M^+\sim(X^-)^\iota(1),
\]
give cyclicity of every rational even component
\cite[Proposition~3.4.7 and Corollary~3.4.8, p.~92]{SharifiIwasawaNotes}.
An even companion root \(\alpha\) of character \(\omega^e\)
corresponds to the target root \(u^2/\alpha\) of character
\(\omega^{2-e}\), agreeing with both the involution and Tate
twists above.  Proposition~\ref{prop:universal-pair} now applies.
\end{proof}

\paragraph{What this removes, and what remains.}
Condition (NC.14) retains the exact rational support condition
of its source; it is not replaced by an untwisted shared-root
test or by finiteness of all \(X^+\).
Greenberg implies these conditions, but no new arithmetic
verification of them with \(X^+[1/p]\ne0\) is asserted.
Combining Corollary~\ref{cor:wwe-pair} with
Theorem~\ref{thm:odd-torsion-cubic} removes Pair-NV as a separate
hypothesis whenever (NC.14) holds.  Rational cyclicity of each
odd torsion character and the actual Triple-NV condition must
still be supplied.  In particular the norm-family argument
does not establish arithmetic cubic noncancellation or a new
unconditional kernel range.
